\section{Implementation}

L'implementation vise deux objectifs : rester proche du modele theorique et offrir une base testable. Les choix techniques sont donc simples, largement disponibles et adaptes a une execution locale.

\subsection{Stack technologique}

Le projet est implemente en Python. Le fichier \texttt{pyproject.toml} declare une compatibilite \texttt{Python >=3.10}. Les principales dependances sont \texttt{pandas} pour la manipulation des tables MovieLens, \texttt{NumPy} pour les calculs numeriques, \texttt{PyYAML} pour la configuration, \texttt{Tkinter} pour la GUI, \texttt{Matplotlib} pour les figures et \texttt{pytest} pour les tests. Le paquet \texttt{click} fournit l'interface CLI.

\subsection{Organisation du code}

L'arborescence du dossier \texttt{src/} suit les modules presentes dans l'architecture :
\begin{verbatim}
src/
  data_manager/    chargement, validation, preprocessing
  fuzzy/           ensembles flous, variables, Mamdani, defuzzification
  recommender/     profils, prefiltrage, scoring, explications
  evaluation/      metriques et protocole train/test
  ui/              commandes CLI et interface Tkinter
  visualization/   generation des courbes d'appartenance
\end{verbatim}

Cette organisation permet de tester separement les composants. Par exemple, la defuzzification peut etre verifiee sans lancer l'interface graphique, et l'evaluation peut reutiliser un moteur deja construit.

\subsection{Configuration externalisee}

Le fichier \texttt{config/fuzzy\_config.yaml} contient les variables linguistiques, les fonctions d'appartenance, les regles, la methode de defuzzification, le seuil de pre-filtrage et les chemins de donnees. Cette externalisation ameliore la reproductibilite : les valeurs du rapport peuvent etre relancees avec le meme fichier de configuration. Elle facilite aussi les ajustements : ajouter des regles ou modifier une fonction d'appartenance ne demande pas de recompilation ni de modification des interfaces.

\subsection{Execution}

Les commandes principales sont les suivantes :
\begin{verbatim}
python main.py recommend --user-id 1 --top-n 10 \
  --set-genre "Sci-Fi=forte,Action=moyenne" --explain

python main.py gui --data-dir data/movie

python main.py evaluate --user-id 1 --top-n 10
\end{verbatim}

La premiere commande produit un Top-N explique, la deuxieme lance l'interface graphique, et la troisieme calcule les metriques d'evaluation. Ces commandes constituent le point d'entree des experiences de la section 9.

La section suivante presente les interfaces qui rendent ces operations accessibles a l'utilisateur.
